% Zenodo DOIs of the code archive. Release v1.0.0 was archived under
% 10.5281/zenodo.22950276 and release v2.0.0 under 10.5281/zenodo.23038895;
% both are versions of the concept record 10.5281/zenodo.22950275.
% \zenodoCodeID is the DOI the paper cites for the code. Until the version
% DOI of release v3.0.0 is minted it is the concept DOI, which resolves to
% the latest version; replace it by the version DOI once Zenodo assigns it.
\newcommand{\zenodoFirstDOI}{%
  \href{https://doi.org/10.5281/zenodo.22950276}{\texttt{10.5281/zenodo.22950276}}}
\newcommand{\zenodoVersionDOI}{%
  \href{https://doi.org/10.5281/zenodo.23038895}{\texttt{10.5281/zenodo.23038895}}}
\newcommand{\zenodoConceptDOI}{%
  \href{https://doi.org/10.5281/zenodo.22950275}{\texttt{10.5281/zenodo.22950275}}}
\newcommand{\zenodoCodeID}{10.5281/zenodo.22950275}
\newcommand{\zenodoCodeDOI}{%
  \href{https://doi.org/\zenodoCodeID}{\texttt{\zenodoCodeID}}}

\section*{Declarations}

\subsection*{Funding}
The authors received no funding for this work.

\subsection*{Competing interests}
The authors have no competing interests to declare.

\subsection*{Ethics approval and consent}
This article reports no research involving human participants, human data or
animals, so ethics approval and consent to participate or to publish were not
required. The results are original and have not been published elsewhere,
and the work of others is credited by citation throughout.

\subsection*{Data availability}
The programs that carry out the computations cited in this paper as
\cite{BMB26}, the Lean~4 certificate of the finite arithmetic, the Macaulay2
computations of local $\Ext$ algebras with their output, the sources of the
paper, and the generators and TikZ sources of every figure are openly
available on Zenodo under the DOI \zenodoCodeDOI. The version that
accompanies this paper is release 3.0.0 of that archive. A document in the
archive lists, for each result of the paper, the computations that check
it. The programs perform $2004$ checks, all of which pass, and the
Lean kernel checks one hundred and ten theorems, each depending on no axiom
or on propositional extensionality alone. The code and the computations were
written and run with the assistance of Claude Code, an AI model developed by
Anthropic. Every result they certify can be reproduced from the archive
with the programs and the Lean kernel, and the authors are responsible for
the mathematical statements and their proofs. No other data were generated
or analysed.
